Buena parte del poder de la teoría de categorías se deriva de su reflexividad:
las categorías son estructuras algebraicas y, como tales, es posible estudiarlas
usando teoría de categorías. Para ello primero debemos ver cuáles son los
morfismos entre categorías, los llamados \emph{funtores}. Este capítulo se
basa principalmente en \cite[pp. 21--32]{joyofcats} y \cite[pp. 13--15]{maclane}.

\begin{definition}
  Un \conc{funtor} entre dos categorías $\cC$ y $\cD$ es un par de funciones
  $T=(o:\Ob{\cC}\to\Ob{\cD},m:\Mor{\cC}\to\Mor{\cD})$ que preserva el dominio, el
  codominio, las identidades y la composición, es decir, tal que:
  \begin{enumerate}
  \item Para cada morfismo $f:a\to b$ en $\cC$, $mf:oa\to ob$ en $\cD$.
  \item Para $f:a\to b$ y $g:b\to c$ en $\cC$, $m(g\circ f)=mg\circ mf$.
  \item Para cada objeto $c$ de $\cC$, $m(1_a)=1_{oa}$.
  \end{enumerate}
\end{definition}

Nótese que la última condición determina unívocamente la función sobre los
objetos a partir de la función sobre los morfismos, por lo que la primera es
redundante. Si $\cC$ y $\cD$ son categorías, escribimos $T:\cC\to\cD$ para decir
que $T$ es un funtor de $\cC$ a $\cD$, y usamos $T$ para referirnos
indistintamente al funtor, a la función sobre los objetos y a la función sobre
los morfismos.

\begin{example}\;
  \begin{enumerate}
  \item Toda categoría $\cC$ admite un \conc{funtor identidad} $1_\cC:\cC\to\cC$
    que asocia a cada objeto o morfismo el propio objeto o morfismo.
  \item Dadas dos categorías $\cC$ y $\cD$ y $d\in\Ob{\cD}$, existe un
    \conc{funtor constante} $C_d:\cC\to\cD$ que lleva todos los morfismos a $1_d$.
  \item La operación <<conjunto potencia>> es un funtor $\power:\bSet\to\bSet$, que
    lleva cada función $f:A\to B$ a la función $(\power f)(S)\coloneqq f[S]$,
    que asocia a cada subconjunto de $A$ su \emph{imagen} por $f$.
  \item De forma similar podemos definir el funtor \conc{conjunto potencia
      contravariante}, $\copower:\bSet\to\dual{\bSet}$, que lleva un conjunto
    $S$ a su potencia $\power{S}$ y una función $f:A\to B$ a la función
    $(\copower f)(T)\coloneqq f^{-1}[T]$ que asocia a cada subconjunto de $B$ su
    \emph{preimagen} por $f$.
  \item Para $n\in\sNat$, existe un funtor $\text{GL}_n:\bCRng\to\bGrp$ que a
    cada anillo conmutativo $C$ le asocia el grupo multiplicativo
    $\text{GL}_n(C)$ de matrices regulares $n\times n$ con entradas en $C$. Los
    homomorfismos de anillos se transforman en homomorfismos de grupos que
    actúan componente a componente.
  \item Sea $\bTop_*$ el constructo cuyos objetos son pares $(X,x)$ formados por
    un espacio topológico $X$ y un punto destacado $x\in X$ y cuyos morfismos
    son funciones continuas que conservan el punto destacado. Entonces podemos
    definir un funtor grupo de homotopía $\pi:\bTop_*\to\bGrp$ que a cada
    espacio topológico $X$ y cada punto $x\in X$ le asocia el grupo de homotopía
    y a cada morfismo en $\bTop_*$ le asocia el correspondiente morfismo de
    grupos.\cite[p. 13]{maclane}

    \begin{proof}
      Primero vemos que la operación sobre morfismos está bien definida. Sea
      $f:(X,x)\to(Y,y)$ un morfismo en $\bTop_*$, y sean
      $\gamma,\sigma:[0,1]\to X$ representantes de un mismo elemento
      $\overline{\gamma}=\overline{\sigma}$ de $\pi(X,x)$, entonces existe una
      homotopía $s:[0,1]\times[0,1]\to X$ de $\gamma$ a $\sigma$, con lo que
      $f\circ s:[0,1]\times[0,1]\to Y$ es una homotopía de $f(\sigma)$ a
      $f(\gamma)$ y $(\pi f)(\overline{\gamma})=(\pi f)(\overline{\sigma})$.
      Además $\pi f$ es un morfismo de grupos, pues la composición con $f$ lleva
      la curva constante en $x$ a la curva constante en $y$ y respeta la
      concatenación de curvas. Finalmente, dados morfismos $f:(X,x)\to (Y,y)$ y
      $g:(Y,y)\to(Z,z)$ en $\bTop_*$, es fácil ver que
      $\pi(g\circ f)=\pi g\circ\pi f$ y que $\pi(1_{(X,x)})=1_{\pi(X,x)}$.
    \end{proof}
  \item Los funtores se pueden componer. Dados dos funtores $S:\cB\to\cC$ y
    $T:\cC\to\cD$, el \conc{funtor composición} $T\circ S:\cB\to\cC$ viene dado
    sobre los objetos como $a\mapsto S(Ta)$ y sobre los morfismos como
    $f\mapsto S(Tf)$.
  \end{enumerate}
\end{example}

\section{Categorías de categorías}

En vista de los ejemplos anteriores sería razonable considerar una <<categoría
de las categorías>>, pero esto plantea ciertos problemas. En primer lugar, esta
categoría no se podría contener a sí misma por la paradoja de Russell. De hecho,
sólo las categorías que son conjuntos pueden estar dentro de una categoría de
categorías, pues las clases propias no pueden estar dentro de otras clases.

\begin{definition}
  Una categoría es \conc{pequeña} si es un conjunto, es decir, si tanto su
  clase de objetos como su clase de morfismos son conjuntos. Llamamos $\bCat$
  a la categoría de las categorías pequeñas y los funtores entre ellas.
\end{definition}

Esto no es del todo satisfactorio, pues la mayoría de las categorías no son
pequeñas. Es por ello que en teoría de categorías es común usar extensiones de
ZFC para lidiar con estos casos. MacLane\cite[pp. 21--26]{maclane} propone una
extensión basada en universos de Grothendieck.

\begin{definition}
  Un \conc{universo} (\conc{de Grothendieck}) es un conjunto $\UNIVERSE$ tal que:
  \begin{enumerate}
  \item Si $x\in u\in\UNIVERSE$ entonces $x\in\UNIVERSE$.
  \item Si $u,v\in\UNIVERSE$ entonces $\{u,v\}\in\UNIVERSE$.
  \item Si $x\in\UNIVERSE$ entonces $\power{x},\bigcup{x}\in\UNIVERSE$.
  \item Si $I\in\UNIVERSE$ y $f:I\to\UNIVERSE$ es una función, entonces
    $\Img{f}\in\UNIVERSE$.
  \item $\sNat\in\UNIVERSE$.
  \end{enumerate}
\end{definition}

La idea es que un universo contendría todos los conjuntos con lo que uno
trataría trabajar normalmente dentro de ZFC, como ilustran las siguientes
propiedades fáciles de probar.

\begin{proposition}
  Sea $\UNIVERSE$ un universo:
  \begin{enumerate}
  \item Si $v\subseteq u\in\UNIVERSE$ entonces $v\in\UNIVERSE$.
    \item Si $u,v\in\UNIVERSE$, entonces $(u,v),u\times v\in\UNIVERSE$.
    \item Si $\{x_i\}_{i\in I}$ es una familia de elementos de $\UNIVERSE$ con $I\in\UNIVERSE$,
      entonces $\prod_{i\in I}x_i,\bigcup_{i\in I}x_i,\bigcap_{i\in I}x_i\in\UNIVERSE$.
    \item Si $a,b\in\UNIVERSE$, todas las funciones $f:a\to b$ cumplen $f\in\UNIVERSE$.
    \item Si $a\in\UNIVERSE$ y $b\subseteq\UNIVERSE$ con $|a|=|b|$ entonces $b\in\UNIVERSE$.
  \end{enumerate}
\end{proposition}

% TODO

%%% Local Variables:
%%% mode: latex
%%% TeX-master: "main"
%%% End:
